\section{مواصفات المنصة: كيف نعيد بناءها}
\label{sec:dishspec}

جمعنا أدناه كل ما يلزم لبناء منصة كهذه من جديد: العتاد، والحلقة، وترميز الدخل، وقراءة الخرج، والتغذية الراجعة، وبروتوكول التجربة. وكل معامل موسوم بمصدره. ”المقالة“: القيمة مأخوذة من نص الدراسة~\cite{Kagan2022}. ”اختيار“: القيمة غير موجودة في المقالة، وعلى من يعيد التجربة أن يحددها بنفسه؛ ونقترح قيمة افتراضية ونبيّن كيف تُختبر.

\subsection*{العتاد والمزرعة}

\begin{center}
\footnotesize
\setlength{\tabcolsep}{4pt}
\renewcommand{\arraystretch}{1.15}
\begin{tabular}{>{\raggedright}p{4.0cm}>{\raggedright}p{6.9cm}>{\raggedright\arraybackslash}p{1.6cm}}
\toprule
المعامل & القيمة & المصدر \\
\midrule
الشريحة & مصفوفة \foreignlanguage{english}{MaxOne}: $26\,000$ قطب من البلاتين على مساحة $8\,\text{mm}^2$ & المقالة \\
التسجيل & حتى $1024$ قناة في آن واحد، $20\,000$ عينة في الثانية & المقالة \\
التحفيز & حتى $32$ قطبًا تقنيًا، وفي التجربة $8$ مستقلة & المقالة \\
الخلايا & قشرة جنين الفأر؛ عصبونات بشرية من خلايا جذعية (طريقتا استنبات)؛ خلايا \foreignlanguage{english}{HEK293T} (ليست عصبونات) للمقارنة الضابطة & المقالة \\
الزرع & نحو $10^6$ خلية للشريحة & المقالة \\
عمر المزرعة عند التجربة & الفأر: من $\sim10$ أيام؛ الإنسان: $14\text{--}24$ يومًا أو $\sim80$ يومًا بحسب طريقة الاستنبات & المقالة \\
\bottomrule
\end{tabular}
\end{center}

\subsection*{الحلقة والزمن}

تعمل الحلقة بنوافذ مدة كل منها $10\ms$ ($200$ عينة): في كل نافذة تُعدّ النبضات على أقطاب المنطقتين الحركيتين، وتُحدَّث اللعبة، ويُصدر التحفيز التالي. والتأخير من النبضة في المزرعة إلى المنبه المجيب نحو $5\ms$. وفي أثناء اللعب تمرّ إشارة كل قطب بمرشح بِسل تمرير عالٍ من الرتبة الثانية بتردد قطع $100\,\text{Hz}$؛ وتُنعَّم القيمة المطلقة للإشارة بمرشح بِسل تمرير منخفض من الرتبة الأولى بتردد $1\,\text{Hz}$، وعتبة النبضة متناسبة مع هذه القيمة المنعَّمة. أما عتبة ستة انحرافات معيارية للخلفية فتذكرها المقالة لقياسات منفصلة للفعالية، لا للعب. وكي لا يُحسب التحفيز استجابةً من المزرعة، تُخمد كل الإشارات حين يصل في آن واحد أكثر من $15$ نبضة كبيرة، تزيد على $75\mV$. كل هذا من المقالة.

\subsection*{ترميز الكرة}

\begin{center}
\footnotesize
\setlength{\tabcolsep}{4pt}
\renewcommand{\arraystretch}{1.15}
\begin{tabular}{>{\raggedright}p{3.4cm}>{\raggedright}p{7.5cm}>{\raggedright\arraybackslash}p{1.6cm}}
\toprule
المعامل & القيمة & المصدر \\
\midrule
المنطقة الحسية & $8$ أقطاب مرتبة بحيث تعني الأقطاب المتجاورة مواضع متجاورة للكرة & المقالة \\
ترميز المكان & رقم القطب $k=1,\dots,8$ يحدد موضع الكرة بالنسبة إلى مركز المضرب & المقالة \\
أي إحداثي & الموضع الرأسي للكرة؛ وللإعادة: بالنسبة إلى المضرب، $k=\lceil 8(y-p+\tfrac12)\rceil$ عند $y-p\in[-\tfrac12,\tfrac12]$ & المقالة؛ والصيغة اختيار \\
ترميز التردد & تردد التحفيز من $4$ إلى $40$ نبضة/ث يحمل المسافة إلى المضرب & المقالة \\
اتجاه السُّلَّم & $4$ نبضات/ث حين تكون الكرة عند الجدار المقابل، وخطيًا حتى $40$ نبضة/ث حين تكون عند جدار المضرب: $f=4+36\,(1-x/L)$، حيث $x$ المسافة إلى جدار المضرب و$L$ عرض الملعب & المقالة \\
شكل النبضة & مستطيلة قصيرة ثنائية الطور: الطور الموجب أولًا، ثم السالب & المقالة \\
السعة & $75\mV$؛ اختيرت كي لا تُجبر الخلايا المثبَّطة على الإطلاق & المقالة \\
مدة الطورين & غير مذكورة؛ تؤخذ الإعدادات القياسية لنظام التحفيز ولا تُغيَّر في أثناء التجربة & اختيار \\
\bottomrule
\end{tabular}
\end{center}

\subsection*{قراءة المضرب}

في المقالة منطقتان حركيتان: النبضات في الأولى تحرّك المضرب إلى أعلى، وفي الثانية إلى أسفل؛ وتُوزن مساهمة المنطقتين لمراعاة اختلاف كثافة الخلايا واختلاف الفعالية~\cite{Kagan2022}. والصيغة الدقيقة غير مذكورة. وللإعادة نأخذ القاعدة~\eqref{eq:dishreadout}، $\Delta p=c\,(n_\uparrow-n_\downarrow)$ لكل نافذة $10\ms$، بوزن يسوّي بين المنطقتين بحسب متوسط فعاليتهما في الراحة (اختيار): يُقسم $n_\uparrow$ و$n_\downarrow$ على متوسط عدّ منطقتيهما في النافذة، مقيسًا قبل اللعب. ونختار المعامل $c$ بحيث يزيح فرقٌ بمقدار عدّ متوسط واحد المضربَ بنسبة $1\%$ من ارتفاع الملعب في النافذة (اختيار)؛ وعندئذ يقطع المضرب الملعب كله في $1\,\text{s}$ إذا رجحت إحدى المنطقتين رجحانًا ثابتًا.

ومواقع المنطقتين على الشريحة لم تُحدَّد في نص المقالة بإحداثيات؛ ليس هناك إلا مخطط. وللإعادة تكفي ثلاث مجموعات غير متداخلة من الأقطاب: ثمانية أقطاب حسية في الوسط، ومجموعة أقطاب تسجيل على كل جانب، إحداهما ”أعلى“ والأخرى ”أسفل“ (اختيار).

\subsection*{اللعبة}

أبعاد الملعب وطول المضرب وسرعة الكرة غير مذكورة في المقالة؛ لا يُذكر إلا أن الكرة تطير بسرعة ثابتة وترتد عن الجدران. وللإعادة (اختيار): ملعب $1\times1$، ومضرب على الجدار الأيمن طوله $0.2$ (إصابة عند $|y-p|<0.1$، كما في النموذج~\foreignlanguage{english}{\texttt{models/dishbrain\_model.py}})، والكرة تقطع الملعب في $1\,\text{s}$. والإخفاق من غير أي كرة مصدودة في الجولة يُعدّ ”إرسالًا ساحقًا“ (\foreignlanguage{english}{ace})؛ والجولة التي تزيد على ثلاث إصابات متتالية تُعدّ ”طويلة“ (المقالة).

\subsection*{التغذية الراجعة}

\begin{center}
\footnotesize
\setlength{\tabcolsep}{4pt}
\renewcommand{\arraystretch}{1.15}
\begin{tabular}{>{\raggedright}p{2.6cm}>{\raggedright}p{8.3cm}>{\raggedright\arraybackslash}p{1.6cm}}
\toprule
الحدث & ما يُقدَّم & المصدر \\
\midrule
إصابة & منبه متوقَّع: الأقطاب الحسية الثمانية كلها في آن واحد، $75\mV$، $100$ نبضة/ث، $100\ms$؛ ويُقطع تحفيز الكرة المعتاد خلال ذلك & المقالة \\
إخفاق & منبه غير متوقَّع: $150\mV$، $5$ نبضات/ث، $4\,\text{s}$، في لحظات عشوائية وعلى أقطاب مختارة عشوائيًا من الثمانية؛ ثم $4\,\text{s}$ بلا تحفيز، وتنطلق الكرة في اتجاه عشوائي & المقالة \\
قانون العشوائية & غير مذكور؛ وللإعادة: لحظات النبضات تدفق بواسوني بمتوسط تردد $5$ نبضات/ث & اختيار \\
\bottomrule
\end{tabular}
\end{center}

\subsection*{البروتوكول والتجارب الضابطة}

تدوم الجلسة الواحدة $20$ دقيقة؛ وتُقارن أول $5$ دقائق بآخر $15$ (المقالة). ومقاييس النتيجة ثلاثة: متوسط طول الجولة، ونسبة الإرسالات الساحقة، ونسبة الجولات الطويلة. وشروط التجربة (المقالة):
\begin{itemize}
\item \emph{المنبه}: الحلقة الكاملة، كما وُصفت أعلاه؛
\item \emph{الصمت}: بدل منبه الإصابة أو الإخفاق المدة نفسها بلا أي تحفيز، ثم تنطلق الكرة في اتجاه عشوائي؛
\item \emph{بلا تغذية راجعة}: بعد الإخفاق ترتد الكرة فحسب وتواصل طيرانها، ويستمر تحفيز موضع الكرة؛
\item \emph{الراحة}: اللعبة تجري والمضرب يتحرك بحسب الفعالية، لكن لا تحفيز أصلًا؛
\item \emph{الضوابط}: شريحة فيها وسط الاستنبات وحده؛ خلايا \foreignlanguage{english}{HEK293T}؛ ”مزرعة في الحاسوب“، أي محاكاة بدل الخلايا.
\end{itemize}
حجم العينة في السلسلة الرئيسية: $9$ مزارع من الفأر ($101$ جلسة)، و$11$ مزرعة بشرية ($138$)، و$6$ شرائح بوسط الاستنبات ($80$)، و$20$ مزرعة في الراحة ($42$)، و$3$ محاكيات ($38$)، والمجموع $399$ جلسة. وفي سلسلة التغذية الراجعة $486$ جلسة على مدى $3$ أيام بمعدل $3$ جلسات يوميًا، بالشروط ”المنبه“ و”الصمت“ و”بلا تغذية راجعة“ ($n=27$ و$17$ و$15$). ولم تظهر زيادة متوسط طول الجولة من أول $5$ دقائق إلى آخر $15$ إلا في المزارع الحية. وفي سلسلة التغذية الراجعة لم يظهر نمو ذو دلالة مع الزمن إلا في شرط ”المنبه“؛ وفي آخر الجلسة تفوّق شرط ”الصمت“ مع ذلك على ”بلا تغذية راجعة“، لكن بأثر أصغر، ولم يكن هناك تعلّم ثابت من جلسة إلى أخرى على مدى ثلاثة أيام~\cite{Kagan2022}.

\subsection*{دورة المنصة خطوةً خطوة}

كل $10\ms$ يفعل حاسوب المنصة ما يلي.
\begin{enumerate}
\item يقرأ عدد النبضات $n_\uparrow$ و$n_\downarrow$ في المنطقتين الحركيتين خلال النافذة المنقضية، ويقسمهما على متوسط عدّ المنطقة في الراحة.
\item يزيح المضرب بمقدار $\Delta p=c\,(n_\uparrow-n_\downarrow)$ ويحصره بين حدّي الملعب.
\item يحرّك الكرة خطوة واحدة؛ ويعكسها عن الجدران العلوي والسفلي والأيسر.
\item إذا بلغت الكرة الجدار الأيمن: عند $|y-p|<0.1$ فهي إصابة، ويزيد عدّاد الجولة واحدًا، ويُقدَّم منبه الإصابة ($100\ms$)؛ وإلا فهي إخفاق، وتُسجَّل الجولة، ويُقدَّم منبه الإخفاق ($4\,\text{s}$)، ثم توقف $4\,\text{s}$، وتنطلق الكرة من جديد.
\item وإلا يختار القطب $k$ بحسب الموضع الرأسي للكرة والتردد $f$ بحسب المسافة إلى جدار المضرب، ويقدّم إلى القطب $k$ نبضات ثنائية الطور بسعة $75\mV$ بالتردد $f$ حتى النافذة التالية.
\end{enumerate}
وهذا يكفي لبناء منصة متوافقة مع المنشورة، ولاختبار السؤال الرئيسي في الفصل: هل يعلّم المزرعةَ التوقُّعُ، أم مجرد الفرق بين الاستجابة للإصابة والاستجابة للإخفاق؟ ولهذا يحسن أن يُضاف إلى شروط المقالة الثلاثة شرط رابع ليس فيها: بعد الإخفاق منبه \emph{متوقَّع} بالمدة والطاقة نفسيهما اللتين للمنبه غير المتوقَّع. فإن تعلمت المزرعة مع ذلك، فالأمر في الفرق لا في التوقُّع.
